\section{个人思考}

\subsection{最重要的理解}

FORGE 的核心价值不在于提出新模型，而在于把制造场景中的“认知难点”拆成可评测任务。它让我们看到，通用 MLLM 在制造任务中的失败往往不是因为完全看不见物体，而是因为缺少型号级语义、零件功能知识、装配规则和微观形态理解。

这一点对后续研究很重要。如果只做更强的视觉 grounding，可能无法解决 FORGE 暴露出的主要问题。模型需要的是制造领域知识注入、细粒度形态建模和任务级推理链训练。

\subsection{对 Benchmark 设计的启发}

FORGE 说明一个好的 benchmark 不能只堆数据量，而要明确任务是否贴近真实决策。制造现场真正关心的是零件是否正确、型号是否匹配、表面是否合格、装配是否满足规格。这些都不是普通分类任务能覆盖的。

此外，FORGE 的任务设计保留了 MLLM 的交互形式：用多项选择题、坐标/字母选项、参考图和上下文示例模拟真实问答。这比只训练一个分类头更接近通用 MLLM 的实际使用方式。

\subsection{三视图策略的利弊}

三视图是一个务实选择。它不等于完整 3D 理解，但能让没有 3D encoder 的通用 MLLM 读取点云的关键几何结构。它保留了工件轮廓、长度、厚度和部分缺陷形态，同时避免坐标序列化带来的长文本和空间结构丢失问题。

但三视图也有天然局限：遮挡、曲率、内部结构和真实三维拓扑仍可能丢失。尤其是 SURFINSP 这类依赖局部表面形态的任务，三视图可能不足以表达微小缺陷。因此，如果后续目标是提升实际工业质检能力，可能需要把三视图和专门 3D/局部几何特征结合起来。

\subsection{复现与扩展建议}

\begin{enumerate}
    \item 先复现三类任务的数据构造和 MCQ 评测协议，确保选项解析、坐标标注和 exact-match 评分一致。
    \item 单独复现 visual grounding probing，区分模型到底是定位失败还是领域推理失败。
    \item 对 SURFINSP 进行更细的错误分析，检查失败来自缺陷不可见、视图信息不足，还是模型缺少缺陷类型知识。
    \item 尝试对小模型做 SFT 时，应采用场景级 held-out 划分，避免只记忆具体零件布局。
    \item 若扩展到新制造场景，应优先建立型号级标注和装配规则标注，而不是只增加粗粒度类别标签。
\end{enumerate}

\subsection{后续问题}

\begin{itemize}
    \item 三视图是否足以支持所有制造几何任务，哪些任务必须引入原生 3D 编码器？
    \item FORGE 的 SFT 增益主要来自学习制造术语、学习推理格式，还是学习具体零件形态？
    \item ICD 在图像模态中提升明显，但在三视图中可能下降，如何设计更适合 3D/多视图的 in-context 示例？
    \item flat washer 等细粒度变体失败是否能通过局部放大图、显式尺寸标注或 CAD 先验缓解？
    \item 如果将 FORGE 用于真实产线，如何处理照明、遮挡、磨损、非标准摆放和相机标定误差？
\end{itemize}
